\section{34 城大气污染防治方案取证表}\label{app:documents}

\noindent 符号：$\checkmark$ 是；$\times$ 否（官方全文归一化检索 0 命中）；? 文件未找到；$\triangle$ 无法判断。\\
证据档：A 全文核；B 官方发布稿／新闻稿；C 第三方镜像或转载；D 未获取。

\begingroup
\footnotesize
\setlength{\tabcolsep}{3.5pt}
\begin{longtable}{@{}R{0.55cm}L{1.15cm}L{1.35cm}C{0.85cm}L{4.35cm}L{4.45cm}C{1.2cm}@{}}
\caption{34 城大气污染防治方案中“烟花爆竹”的出现情况}\label{tab:34-app}\\
\toprule
\textbf{\#} & \textbf{城市} & \textbf{文件年} & \textbf{提到} & \textbf{独立列项} & \textbf{牵头单位（原文表述）} & \textbf{档} \\
\midrule
\endfirsthead
\toprule
\textbf{\#} & \textbf{城市} & \textbf{文件年} & \textbf{提到} & \textbf{独立列项} & \textbf{牵头单位（原文表述）} & \textbf{档} \\
\midrule
\endhead
\bottomrule
\endfoot
1 & 东营 & 2023 & $\checkmark$ & \textbf{7 个专项行动之第（六）项，下辖第 18 项} & 市公安局、市应急局、市司法局、市生态环境局 & A \\
2 & 济南 & 2023\slash2024 & $\checkmark$ & 面源治理第（四）条；2024 版为非独立的一款 & 市公安局、市应急局 & A \\
3 & 青岛 & 2024 & $\checkmark$ & \textbf{5 大方面之一、十二项重点工作第 12 项} & 未写明 & B \\
4 & 淄博 & 2023 & $\checkmark$ & 面源治理提升行动第（四）项 & 公安列首但未写“牵头” & A \\
5 & 潍坊 & 2020\slash2024 & $\times$ & — & — & A \\
6 & 郑州 & 2024 & $\checkmark$ & \textbf{8 大方面 35 项之第 25 项} & 市应急局、市公安局 & A \\
7 & 洛阳 & 2024 & $\checkmark$ & 面源治理第（十八）项 & 市应急管理局牵头 & A \\
8 & 南阳 & 2024 & $\checkmark$ & 面源治理第（四）项 & 未写“牵头”，应急管理局列首 & A \\
9 & 周口 & 2020 & $\checkmark$ & \textbf{10 个部分 52 项之第 20 项} & 未写“牵头”，应急管理局列首 & A \\
10 & 安阳 & 2024 & $\checkmark$ & 八大攻坚行动第 6 项下的分项（7 项之末） & 未写明 & B \\
11 & 滨州 & 2024 & ? & — & — & D \\
12 & 佛山 & 2017 & $\checkmark$ & 倡导性；“全面禁止露天焚烧”条的末句 & \textbf{市民族宗教局}（与“文明敬香、生态寺院”合并） & A$^{\dagger}$ \\
13 & 成都 & 2024 & $\checkmark$ & \textbf{9 方面 40 项之第 24 项独立任务} & 公安、生态环境、交通运输、应急、邮政管理 & A$^{\ddagger}$ \\
14 & 昆明 & 2025 & $\times$ & — & — & A \\
15 & 南京 & 2019\slash2024 & $\checkmark$ & 2024 版为非独立末两句；2019 版为\textbf{第 36 项独立任务} & 2024 版未写明；2019 版市公安局牵头 & A \\
16 & 苏州 & 2024 & $\checkmark$ & \textbf{第（十四）条独立成条 $+$ 附件任务表第 27 项} & 公安列首但未写“牵头”（6 部门分工） & A \\
17 & 杭州 & 2024 & $\times$ & — & — & A \\
18 & 宁波 & 2025 & $\checkmark$ & 全文仅一处；“强化执法监管”的末句 & 未写明（全篇 0 次“牵头”） & A \\
19 & 合肥 & 2025 & $\checkmark$ & 联防联控第（二十一）项中的一句 & \textbf{市生态环境局牵头} & C$^{\ddagger}$ \\
20 & 广州 & — & ? & — & — & D \\
21 & 东莞 & 2025 & $\triangle$ & — & — & D$^{\S}$ \\
22 & 绵阳 & — & ? & — & — & D \\
23 & 石家庄 & 2021 & $\checkmark$ & \textbf{60 项之第 52 项} & \textbf{市公安局}（配合应急、市场监管、供销社、生态环境） & C \\
24 & 唐山 & 2024 & $\checkmark$ & 全篇仅一句，非独立，挂在“秸秆综合利用和禁烧”末尾 & 未写明 & A \\
25 & 太原 & 2024 & $\times$ & — & — & A \\
26 & 临汾 & 2024 & $\checkmark$ & 6 项攻坚任务之一项下的 1 条 & 市级未列（原文 404，仅见官方图解） & B \\
27 & 沈阳 & 2025 & $\times$ & — & — & A \\
28 & 鞍山 & 2024 & $\times$ & — & — & A \\
29 & 丹东 & 2024 & $\times$ & — & — & A \\
30 & 大连 & 2024 & $\checkmark$ & 非独立，一句话；“强化油烟、恶臭等群众身边污染治理”内并列的一项 & 未写明（7 部门按职责分工） & A \\
31 & 武汉 & 2022 & $\checkmark$ & 一句话；“加强大气面源治理”下第 3 条内 & 未写明“牵头”；责任单位列名第一为市公安局 & A \\
32 & 宜昌 & 2024 & $\checkmark$ & 10 个专项行动之第（六）项结尾一句；附件\textbf{第 73 条} & \textbf{市生态环境局}（附件注明“排名第一的为牵头单位”） & A \\
33 & 西安 & 2025 & $\checkmark$ & \textbf{8 个专项攻坚行动之第（二）项“烟花爆竹禁售禁放攻坚行动”，含 6 条子任务} & \textbf{市公安局}（子条 1 与市应急管理局并列；子条 3 为市应急管理局） & A \\
34 & 宝鸡 & 2022 & $\checkmark$ & \textbf{30 项之第 25 项} & \textbf{市公安局牵头} & A \\
\end{longtable}

\vspace{0.3em}
\noindent{\scriptsize $^{\dagger}$ 文件已于 2020-05-13 失效。\quad $^{\ddagger}$ 征求意见稿，正式印发版是否保留未核。\quad $^{\S}$ 方案原文未公开，仅有解读稿（0 命中）。}
\endgroup

\vspace{0.6em}
\noindent\textbf{合计}：34 城中，取得原文的 30 城为分母——提到 \textbf{23/30 $\approx$ 77\%}；独立列项 \textbf{11/30 $\approx$ 37\%}；官方全文零命中 7 城（潍坊、太原、沈阳、鞍山、丹东、昆明、杭州）；文件未找到 4 城（滨州、广州、东莞、绵阳）。若把 4 个未找到的城市计入分母并视为“没有”，则提到的比例为 23/34 $\approx$ 68\%。两种口径都列明，正文采用前者。

\begin{table}[H]
\centering
\caption{层级对照：各层级文件中是否含“烟花爆竹”}
\label{tab:levels}
\small
\begin{tabular}{L{8.4cm}L{6.2cm}}
\toprule
\textbf{层级} & \textbf{含“烟花爆竹”？} \\
\midrule
国家 2017（环大气〔2017〕110 号，32 项） & \textbf{含}，第 25 项，未标牵头单位 \\
国家 2023（国发〔2023〕24 号 39 项；环大气〔2023〕73 号） & \textbf{零命中} \\
河南省级（豫政〔2024〕12 号） & \textbf{含} \\
山东省级（鲁政字〔2024〕102 号） & \textbf{不含} \\
\bottomrule
\end{tabular}
\end{table}

\noindent\textbf{四条限定}：①“未找到”不等于“没有”，4 城不计入分母（正文以 30 城为分母）；②牵头单位各地不一，\textbf{不得写成“由公安局牵头”}；③成都、合肥为征求意见稿，正式印发版是否保留未核；④石家庄为 2021 年版且用第三方镜像，临汾市级原文 404，二者单独引用时须注明。
